\section{Appendix: Adaptive-Attacker Robustness}
\label{sec:adaptive-appendix}

\draftnote{DRAFT. Underlying data and decomposition are final
(Checkpoint 3d, \texttt{reports/03d\_gap\_decomposition.md}); this
section restates it under the PI's relabeling, not as a new analysis.}

An earlier phase of this project asked whether a fairness-aware
adaptive adversary -- one trained and evaluated within a user's true
behavioral decile, with oracle knowledge of that decile -- recovers
substantially more risk than the naive attacker of
Section~\ref{sec:threat-model-naive}. We decompose that adversary's
apparent advantage into three components, computed per user and
bootstrapped over users: (a) naive, the plain attacker;
(b) \emph{external side information about regularity class} -- the
\emph{same} trained naive model, with its output logits masked to only
the user's true decile's candidate users before taking top-1 (no
retraining); (c) specialized, a model trained and tested only on the
true decile's users; (e) realistic, the specialized model from (c) but
routed at test time by a classifier's \emph{predicted} decile (trained
on released features alone) rather than the oracle label.

Table~\ref{tab:gap-decomposition} reports, per dataset, the mean across
all ten deciles of each gap component -- side information $(b-a)$,
specialization $(c-b)$, and the realistic gap $(e-a)$ -- together with
how many of the ten (dataset, decile) cells have a bootstrap CI
excluding zero. Side information alone (b$-$a) is positive and
significant in every single cell across all three datasets: nearly all
of the originally reported ``adaptive adversary'' advantage is
explained by restricting the naive model's candidate pool to the
user's true decile, which requires no attacker sophistication at all.
Specialization (c$-$b) -- actually training a dedicated model once the
candidate pool is already restricted -- is at or below zero in the
large majority of cells, and never reliably positive. Most
importantly, the group-aware attack from released features alone adds
\emph{no risk}: the realistic gap $(e-a)$ is negative in the large
majority of cells across all three datasets, meaning a realistic
adaptive adversary -- one without oracle access to a user's true
behavioral group -- performs no better, and typically worse, than the
naive attacker it was meant to improve on.

\begin{table*}[t]
\centering
\caption{Gap decomposition (Checkpoint 3d): mean of the side-information, specialization, and realistic gap components across all deciles, per dataset, with the count of (dataset, decile) cells whose bootstrap CI excludes zero. Generated from \texttt{reports/figures/phase3d/*/summary.json}.}
\label{tab:gap-decomposition}
\small
\begin{tabular}{@{}lccc@{}}
\toprule
Dataset & Side info.\ (b$-$a) & Specialization (c$-$b) & Realistic (e$-$a) \\
\midrule
GeoLife & +0.166 (10/10 CI excl.\ 0) & -0.030 (3/10 CI excl.\ 0) & -0.041 (9/10 CI excl.\ 0) \\
Foursquare & +0.193 (10/10 CI excl.\ 0) & -0.040 (8/10 CI excl.\ 0) & -0.155 (9/10 CI excl.\ 0) \\
Gowalla-NYC & +0.210 (10/10 CI excl.\ 0) & -0.016 (1/10 CI excl.\ 0) & -0.108 (9/10 CI excl.\ 0) \\
\bottomrule
\end{tabular}
\end{table*}


This finding does not undermine Section~\ref{sec:disparate-risk}'s
disparate-risk result, which depends only on the naive attacker and is
unaffected by any property of the adaptive adversary studied here.
